% Part: lambda-calculus
% Chapter: syntax
% Section: abbreviated-syntax

\documentclass[../../../include/open-logic-section]{subfiles}

\begin{document}

\olfileid[hi]{lam}{syn}{abb}
\olsection{संक्षिप्त वाक्यविन्यास}

\olref[trm]{defn:term} में परिभाषित पद कभी-कभी लिखने में बोझिल होते हैं,
इसलिए अधिक संक्षिप्त वाक्यविन्यास प्रस्तुत करना उपयोगी है। निश्चय ही हमें
यह सुनिश्चित करने में सावधानी रखनी होगी कि संक्षिप्त संकेत-लिपि के पद भी
अद्वितीय रूप से पठनीय हों। व्यापक रूप से उपयोग होने वाला एक रूप, जिसे
\emph{संक्षिप्त पद} कहते हैं, इस प्रकार है।

\begin{enumerate}
\item जब कोष्ठक छोड़ दिए जाएँ, तब अनुप्रयोग बाएँ से दाएँ होता है।
  उदाहरणार्थ, यदि $M$, $N$, $P$ और~$Q$ पद हैं, तो $MNPQ$,
  ~$(((MN)P)Q)$ का संक्षिप्त रूप है।
\item फिर, जब कोष्ठक छोड़ दिए जाएँ, तब लैम्ब्डा अमूर्तन को संभव अधिकतम
  व्याप्ति दी जाती है। उदाहरणार्थ, $\lambd[x][MNP]$ को
  $(\lambd[x][MNP])$ पढ़ते हैं।
\item एक लैम्ब्डा से अनेक चरों का अमूर्तन किया जा सकता है। उदाहरणार्थ,
  $\lambd[xyz][M]$, ~$\lambd[x][\lambd[y][\lambd[z][M]]]$ का संक्षिप्त रूप
  है।
\end{enumerate}

उदाहरणार्थ,
\[
\lambd[xy][xxyx \lambd[z][xz]]
\]
इसका संक्षिप्त रूप है:
\[
(\lambd[x][(\lambd[y][((((xx)y)x)(\lambd[z][(xz)]))])]).
\]

\begin{prob}
संक्षिप्त पद $\lambd[g][(\lambd[x][g (x x)])
  \lambd[x][g (x x)]]$ को पूर्ण रूप में फैलाएँ।
\end{prob}

\end{document}
